% ============================================
% DOC-02 (ZH)
% 中文翻译 — SPEC-935-R356
% ============================================
\section*{\textbf{DOC-02} --- 患者1,260病历}
\markboth{DOC-02 --- 患者1,260病历}{}

\noindent \textit{巴尔巴塞纳收容医院 —— 入院登记表}

\noindent \textbf{巴尔巴塞纳收容医院 —— 入院登记表}
\textit{虚构重构的入院登记表；并非真实档案文件。}
\textit{依据历史背景创作，未经米纳斯吉拉斯州公共档案馆认可。}

\vspace{0.5em}\noindent\rule{\textwidth}{0.5pt}\vspace{0.5em}

\noindent \textbf{登记号:} 1,260  
\textbf{姓名:} 若阿金·安东尼奥·科雷亚  
\textbf{出生日期:} 约1907年（估计）  
\textbf{籍贯:} 塞阿拉州基谢拉莫宾  
\textbf{父母:} 父：安东尼奥·科雷亚（已故 —— 饿毙，1915年）；母：玛丽亚·杜罗萨里奥·科雷亚（失踪 —— 被强制重新安置，1915年）  
\textbf{来源:} 1915年从塞纳多尔·蓬佩乌出发的逃荒者，后被关押在福塔莱萨附近的阿拉加迪索（Alagadiço）收容营。1917年经贝洛奥里藏特圣慈悲救济院转送。

\vspace{0.5em}\noindent\rule{\textwidth}{0.5pt}\vspace{0.5em}

\noindent \textbf{入院检查 —— 内尔松·丰特斯医生，1917年11月12日:}

\noindent \textit{“患者呈严重恶病质状态。估计体重：20公斤（同龄预期：28—30公斤）。明显蛋白质—热量营养不良。下肢水肿。腹部膨胀。牙齿受损（多处龋齿，缺失3颗）。}

\noindent \textit{最显著体征：右眼虹膜异色，呈金黄色，左眼无此现象。患者声称是‘一个怪物’改变了颜色。可能为未明确的感染或外伤后遗症。}

\noindent \textit{精神状态：平静，合作，但目光呆滞茫然。回答问题仅用单音节。被问及家人时称：‘都死了。全死了。’无感觉异常迹象。时间与空间定向力正常。}

\noindent \textit{初步评估：青春期精神分裂症（早发性痴呆）。建议隔离观察。”}

\vspace{0.5em}\noindent\rule{\textwidth}{0.5pt}\vspace{0.5em}

\noindent \textbf{事件记录:}

\begin{moltablethree}
\hline
日期 & 事件 & 后果 \\
\hline
1920年2月 & 无端攻击护工 & 隔离30天 \\
1922年6月 & \textbf{杀人 —— 护工瓦格纳} & 隔离90天 \\
1937年3月 & \textbf{杀人 —— C区患者} & 永久隔离 \\
1937年4月 & 自残（用牙齿拔掉指甲） & 身体约束 \\
1945年9月 & 无法控制的发声（声称‘裹尸布撕裂者’正借他之口说话） & 溴化物镇静 \\
1948年1月 & \textbf{杀人 —— 访客（闯入禁区）} & 终身隔离 \\
1949—1956 & 选择性缄默（7年不与工作人员说话） & --- \\
1956年 & 恢复说话：第一句话‘日记’ & --- \\
1979年 & 埃托尔·奥利维拉医生会谈 & 见附件报告 \\
\hline
\end{moltablethree}

\vspace{0.5em}\noindent\rule{\textwidth}{0.5pt}\vspace{0.5em}

\noindent \textbf{用药记录（部分）:}

\noindent - 溴化钾（1920—1945）—— 镇静剂
- 电休克（1937、1945、1952）—— 无麻醉
- 经眶额叶切除术 —— \textit{1948年曾建议但未实施（患者对仅仅试图接近就做出极端暴力反应）}
- 氯丙嗪（1955年起）—— 抗精神病药

\vspace{0.5em}\noindent\rule{\textwidth}{0.5pt}\vspace{0.5em}

\noindent \textbf{最终意见:}

\noindent 患者1,260是收容院住院时间最长的院民之一。其行为不可预测：绝对平静的时期与极端暴力的爆发交替出现。他保存着一本日记，生前从未允许任何工作人员阅读。

\noindent \textit{补充注记（1979年）：埃托尔·奥利维拉医生获准阅读该日记。他是第一个也是唯一一个接触过此文件的人。奥利维拉医生在48小时后申请病假。再未返岗。}

\noindent \textit{该日记目前由巴尔巴塞纳法医研究所保管，附于尸检报告中。}

\vspace{0.5em}
\noindent\textit{\textbf{↪ 路线:} \rota{MEM-08} \rotasep \rota{DOC-03} \rotasep \rota{LUC-01}}
